\section{Conclusions}
\label{md:conclusiones}
The unifying result of this work is the joint determination of structures
that scalar evaluations present separately. Explicit compatibility and
reconstruction maps have been constructed, the invariants of their
transformations have been identified, and the observations that recover
the retained information have been established. The composition relates
the constants through their shared operations.
Action participates in a dimensional scale, in the counterangle, and in
the electronic return quotient. The angular pair determines oriented
channels; hexad--octad incidences produce constitutive readout operators
and the Barbero--Immirzi functional. The composition preserves this
genealogy and establishes joint properties beyond the individual
evaluation of each readout operator. The starting point remains the
APP--TRIT--TPK composition: the inverse relations act on its constructed
outputs and preserve the orientations, domains, and dimensional data
that they use.

\paragraph{Chronological memory and operator response.}
The chronological analysis specifies the information retained by an
extension. In the capacity factor, with the stated phase measure, block
entropy grows without bound while its rate per position tends to zero.
The balance, which takes two values, loses an increasing amount of
chronological information. The quadratic activity of the current readout
operator remains positive even when its mean vanishes. These properties
concern different observations of the same process and can coexist
without identifying the symbolic rate with a temperature.

Operator reconstruction adds a dynamical distinction. In the nonadic
representation of readout and memory, order sensitivity is distributed
between the contributions \(1/81\) and \(8/81\) of the internal
commutator. Reintroducing memory combines them and produces \(1/9\).
The retained history therefore participates in composing the response:
knowledge of the present balance and knowledge of the capacity to respond
to a continuation constitute different information.

\paragraph{Spectral determination and reconstruction stability.}
Spectral reconstruction establishes a joint determination result. In the
normalized constitutive collection, velocity fixes the product of the two
responses, and the Barbero functional distinguishes their oriented
distribution monotonically. Together, these data recover the labelled
eigenvalues. On the generated image, inversion composes with angular
reconstruction and reconstruction of the action coefficient. Sheet
markers and dimensional bases are retained: they are representation data
used in the reconstruction.

The differential characterization makes this reconstruction precise.
The logarithms of velocity and impedance are derivatives of a single
spectral potential, defined through the incidence channels. Its negative
definite Hessian yields an exact reciprocity of sensitivities and bounds
the inverse on controlled domains. Uniqueness for exact data and
sensitivity to finite errors are thus characterized by distinct results.
This distinction is particularly relevant when strong contraction
preserves injectivity while reducing the observable separation between
states.

\paragraph{Constitutive composition and electroweak reconstruction.}
The cubic inverse transports multiplication of contractions to an
associative law on responses. Its logarithmic coordinate is additive
and retains the direction of composition on each sheet. Electroweak
reconstruction specifies another complementarity: the ratio of gauge
couplings recovers the common contraction, whereas Higgs determines
the oriented coordinate within the specified domain. Marked projectors
turn that scalar recovery into reconstruction of the constitutive
operator. This preserves a proved connection between distinct readers
of the same channels.

\paragraph{Catalan's constant and the constitutive potential.}
The depth-two moment connects two evaluations of the same functional
structure. On the oriented quarter-turn, its imaginary part produces the
collection whose extremal value is Catalan's constant; on the contractive
channels, the same moment enters the constitutive potential. The
arguments, orientation, and normalization of each evaluation remain
specified. The functional structure carries more information than any
one of its extremal values considered in isolation.
The integral formula in Proposition~\ref{md:compos:inversa-catalan}
also recovers the moment from the constitutive response. The zero-limit
condition at the origin uniquely fixes the integration constants.
Production and reconstruction are thus linked on a single functional
collection, together with its domain and endpoint value.

\newpage
\paragraph{Reconstruction of action and memory of order.}
Relative composition shows that order leaves a recoverable operator
signature. The quarter-turns of conjugate forms and their inverses
produce a unimodular loop with factor
\[
 \rho_{\rm lazo}=
 \left(\frac{501\alpha+\eta_{\rm ret}}
 {499\alpha-\eta_{\rm ret}}\right)^2,
 \qquad 0<\eta_{\rm ret}<499\alpha.
\]
The stated branch and references allow \(\eta_{\rm ret}\) to be
recovered. Its dimensional scale remains the necessary antecedent for
reconstructing the Planck section. The lift to the twelve-phase cycle
establishes where this transformation acts, and iteration expresses its
accumulation through reciprocal powers. Nonadic holonomy retains its own
domain and update law; the composition studied here constitutes a
well-defined subsequent transport.

Operational information thereby acquires a precise definition: two
histories are distinguishable with respect to the admissible preparations
and detectors when their transports produce different responses.
Operation multiplicities do not generally determine this response;
order can be indispensable. The trace and vector responses likewise
retain different data. A sufficient collection of readouts reconstructs
the operator, whereas other readouts assign the same value to oppositely
oriented transformations. This definition makes the relation between observation,
memory, and reconstruction explicit, while distinguishing a
representation from the complete genealogy of the state.

\paragraph{Mass reciprocity and the interrelation of physical quantities.}
Joint normalization has a prior operator proof: the inscribed return
produces $L_\alpha$, and polar transport of the same character determines
the readout operators $R$ and $\Lambda$.
The identities $H\Lambda=\hbar cI$ and $R\Lambda=L_\alpha^2I$ fix
$G=c^3L_\alpha^2/\hbar$ as the unique coefficient of the radial realization.
The proof retains noncommutative variations and compatible elimination
of memory, including dual transport of the action length.
Its returned evaluation in the SI section is
$6.674299659414022706\ldots\,10^{-11}\,\mathrm{m^3\,kg^{-1}\,s^{-2}}$;
the CODATA comparison follows that evaluation.

The composition brings together four readouts of a single positive
structural mass coordinate. Under the stated circular and thermal
conditions, the proper frequency, reduced gravitational radius, reduced
Compton wavelength, and energy normalized by the tritic decision satisfy
\[
 \Xi_\Gamma=
 \frac{\Omega_\Gamma}{\omega_0}
 =\frac{r_{g,\Gamma}}{R_{108}}
 =\frac{R_{108}}{\bar\lambda_\Gamma}
 =\frac{E_\Gamma}{k_B\Theta_{108}\log3}.
\]
The electronic prefactor, character, and towers belong to the
construction of \(\Xi_\Gamma\). The identity preserves their roles
and normalizations. In this realization, the mass factor is
characterized as a reciprocal dilation parameter for the two lengths;
their product preserves the area \(R_{108}^2\). The thermal weight of
the cycle is expressed by \(3^{-\Xi_\Gamma}\), subject to the
corresponding normalization and convergence conditions.
Here $R_{108}=L_\alpha$; the circular form preserves the length already
constructed. Varying the clock also transports the electronic quotient
$b_e$, preserving the same mass readout. Chronological windows,
reduced phases, and terminal costs describe aspects of the state;
preserving them in isolation allows changes that the full composition
distinguishes.

\paragraph{Electronic scale, Fermi and action length.}
Substitution of the complete electronic section retains its prefactor
and return within the electroweak realisation. The common scale enters
tree-level Fermi with an inverse-square power and the Higgs vacuum
expectation value with a direct power; it cancels in angular ratios and
in the self-coupling. Species-specific factors remain in their ratios.
Positive congruence also transports the Compton length by inverse
congruence, without requiring commutation between basal mass and return.
These proofs distinguish the absolute scale, angular information and
spectral structure preserved by each operation.

\paragraph{Electronic reconstruction and preservation under transport.}
The electronic readout retains a verifiable dependence on order.
The sum and differences between opposite positions reconstruct the
record, including its integrality conditions; autocorrelations and
the spectrum distinguish records with equal multiplicities. The
coefficient $80$ evaluates the specified central record and changes
when its order is altered. For a finite collection of multisections,
the mass operator is positive, and its commutator with a transport
determines exactly when that transport preserves mass. This criterion
transmits the same preservation to the temporal, spatial, gravitational,
and thermal readouts, with their references fixed.

\newpage
\paragraph{Marked trace and thermometric response.}
The composition of capacities and occupations produces
\[
 b_*=10^{-23}+380\,10^{-26}+649\,10^{-29}
     =1.380649\times10^{-23}.
\]
The grading of the $649$ states retains the sectors $1$, $24$,
and $624$ and their respective memory indices. The spectrum with
ratio $1/10$ follows from the memory transport already proved.
An explicit thermometric realization is constructed on that reference:
marked information is $b_*s(\rho)$, and energy is the readout
conditioned on the marked sector. For Gibbs variations with fixed
Hamiltonian and reference, the temperature satisfies
$b_*\Theta=\beta^{-1}$. The coefficient is thereby characterized
through a concrete thermal operation. Conditioning both quantities
in the same way gives a different normalization; varying the reference
adds the corresponding differential term. The proof identifies these
dependencies and retains the dimensional chart required for the
physical Boltzmann readout. Composition with the tritic clock,
entanglement area, and gravitational reciprocity yields the joint
relation among decision energy, action, and area.

\paragraph{Physical content of the transformation.}
The analysis constructively distinguishes passive transport from
relative response. The former telescopes upon return of the chart and
phase. The latter retains a final operator different from the identity
when the commutator is nontrivial. Symplectic area preservation coexists
with the absence of a common positive energy metric for the two
conjugate forms. Interpreting an energy gain requires the realization of
the controller and the reference to be included. The mathematical result
determines what must be tested and which data must be retained to
attribute a response to the physical system.

\paragraph{Structural definition and unity of the result.}
The developments refine the definitions used throughout the work. The
action coefficient characterizes a regional section that participates
both in angular construction and in its reconstruction through
transport. The Barbero--Immirzi functional distinguishes the oriented
distribution of constitutive channels; together with normalized
velocity, it determines their spectrum. Catalan's constant appears as
an extremal evaluation of the oriented moment whose collection also enters
the response potential. In the specified realization, the mass
coordinate parametrizes a reciprocal dilation preserving the product
of the lengths. Boltzmann's constant relates the temperature and
decision-energy lines through the stated positive transducer and tritic
normalization. Each definition supplements the value with an operation,
a domain, and a role within the composition.

\paragraph{Critical separation and reconstruction of readers.}
Parametric scaling is established for the sequence selected by the construction itself: the first return roots of the uniform dodecaphasic reader. APP--TRIT--TPK, the joint continuum, orientation and memory remain its antecedent. Identification of centres and remainder control in \ref{hmtfg:escalado:15}--\ref{hmtfg:escalado:17} culminate in \eqref{hmtfg:articulacion:limite}: the Feigenbaum constant is the asymptotic separation invariant of those parameters. The same development recovers the profile from its twelve-node measure through Jacobi and, through constitutive inversion, recovers the oriented channels and the moment whose endpoint is Catalan's constant. The result thus brings generation, selection and reconstruction together, preserving the domain and information that distinguish each reader.

% Adición a las conclusiones de VII; las conclusiones anteriores se conservan.
\paragraph{Integration of the four interactions and total canonical closure.}
Gravitation and the strong, weak and electromagnetic interactions act on the same material realization of the APP--TRIT--TPK structure. Transport preserves the sheets, orientation, memory and incidence of the enriched state; its spin and gauge realizations compose on the same material fibres. The electromagnetic component belongs to the electroweak realization: it is not introduced as an additional interaction disconnected from it. The total action uses the coefficients and readers already constructed, including the Planck section, radial length and Barbero--Immirzi functional.

The content of this integration is not limited to collecting terms in an expression. Variation of the same action produces the total spin current and the joint metric source. Eliminating contorsion preserves the cross-terms between species; the Legendre transformation preserves material momenta and produces the total constraints. The proofs in sections~\ref{hmtcuatro:gea:seccion} and~\ref{hmtcuatro:accion-retroaccion-restricciones} establish backreaction, the first-class property and propagation of those constraints in the specified regular realization.

In particular, let $\Theta_{\rm tot}$ be the total canonical potential obtained in that transformation, $\mathcal Z=\{C_A=0\}$ the regular constraint surface, and $X_N=N^AX_{C_A}$, $X_M=M^BX_{C_B}$ two admissible characteristic deformations. The connection induced by the action has $H_N^{\rm tot,can}=-\Theta_{\rm tot}(X_N)$. Its curvature, calculated without omitting derivatives of the deformation coefficients, is
\begin{equation}
\begin{aligned}
\mathcal F^{\rm tot,can}_{N,M}
&=\delta_NH_M^{\rm tot,can}-\delta_MH_N^{\rm tot,can}
 -H_{[X_N,X_M]}^{\rm tot,can}
 +\frac{i}{\hbar}[H_N^{\rm tot,can},H_M^{\rm tot,can}]\\
&=-\mathrm d\Theta_{\rm tot}(X_N,X_M)
 =N^AM^B f_{AB}{}^{D}C_D.
\end{aligned}
\label{hmtcuatro:conclusion-curvatura-total}
\end{equation}
Therefore,
\[
 \left.\mathcal F^{\rm tot,can}_{N,M}\right|_{\mathcal Z}=0.
\]
The formula includes the tangential deformation and the gauge actions of the total bracket; on the quotient, it expresses the same equality modulo those vertical directions. The prequantum representation preserves this closure through the commutator identity proved in \eqref{hmtcuatro:representacion-precuantica}.

Thus, gravitational incorporation is part of the same action, its sources and its constraint algebra, not an equation juxtaposed after constructing the other interactions. The proved vanishing concerns the canonical characteristic curvature; it does not require spacetime curvature or gauge-field strengths to vanish. Its domain preserves the nondegenerate coframe, the regular chart and the boundary conditions used in the proof. Nor does it suppress the memory or global monodromy of the histories that gave rise to the realization.


Horizon composition distinguishes the action dependence at fixed area
from that at fixed mass. Section~\ref{ins29:vii:horizonte} also proves
a compatibility curve connecting thermal bias, elliptical shape, and
the Barbero channels. Probability ratios recover the same action section
and provide a joint check of these readings.


Production, compatibility, and reconstruction therefore constitute three
complementary levels of a single description. Genealogy determines the
objects, joint relations constrain their readouts, and detectors
determine the recoverable structure. The new characterizations rest on
the operations presented here and retain their conditions, allowing
them to be incorporated into specialized developments while preserving
the unity of the argument.
